% =====================================================
% 题型：椭圆方程与直线相交三角形面积及夹角正切最值解答题 (q_ellipse_pqr_tangent.tex)
% 知识板块：椭圆综合
% 主思维：数形结合与几何直观
% 副思维：最值、分类、边界、卷面表达
% 错因标签：斜率/韦达/象限漏条件、算力崩溃
% 讲义用途：解析几何综合
% =====================================================
% 难度：★★★ (难题)
% =====================================================
\newcommand{\qEllipsePqrTangentStem}{%
  已知椭圆 \(C: \dfrac{x^2}{a^2} + \dfrac{y^2}{b^2} = 1 \ (a > b > 0)\) 的左焦点为 \(F(-1, 0)\)，离心率为 \(\dfrac{1}{2}\)。\\
  (1) 求 \(C\) 的方程；\\
  (2) 设 \(O\) 为坐标原点，过 \(F\) 且斜率大于 0 的动直线 \(l\) 与 \(C\) 交于 \(P\)，\(Q\) 两点，其中 \(Q\) 在第三象限，直线 \(PO\) 与 \(C\) 的另一个交点为 \(R\)。\\
  (i) 若 \(\triangle PQR\) 的面积是 \(\triangle PFO\) 的面积的 3 倍，求 \(l\) 的方程；\\
  (ii) 求 \(\tan \angle PQR\) 的最小值。%
}

\newcommand{\qEllipsePqrTangentAnswer}{%
  (1) \(\dfrac{x^2}{4} + \dfrac{y^2}{3} = 1\)； (2)(i) \(y = \dfrac{\sqrt{5}}{2}(x + 1)\)； (ii) 最小值为 \(4\sqrt{3}\)。%
}

\newcommand{\qEllipsePqrTangentSolution}{%
  【核心思维】\\
  第一步，利用椭圆的左焦点和离心率求出参数 \(a, b, c\)，从而确定椭圆的标准方程；第二步，在第二问中，利用极坐标法（焦点弦参数化）来表示相交弦。设直线 \(l\) 与 \(x\) 轴正向的夹角为 \(\theta\)，将点 \(P, Q\) 用焦点极半径表示，由对称性 \(R = -P\)，利用行列式（或向量叉乘）计算出 \(\triangle PQR\) 与 \(\triangle PFO\) 的面积表达式，通过比值求得倾角的余弦值 \(u = \cos\theta\)，进而写出直线方程；第三步，利用向量的夹角正切公式（叉乘除以点乘）表示出 \(\tan\angle PQR\) 关于 \(u\) 的无理代数式，通过二次换元法，转化为求导最值问题。\\
  【标准作答与步骤分析】\\
  解：(1) 因为左焦点为 \(F(-1, 0)\)，所以 \(c = 1\)。\\
  由离心率 \(e = \dfrac{c}{a} = \dfrac{1}{2}\) 得 \(a = 2\)。\\
  所以 \(b^2 = a^2 - c^2 = 4 - 1 = 3\)。\\
  椭圆 \(C\) 的标准方程为：\(\dfrac{x^2}{4} + \dfrac{y^2}{3} = 1\)。\\
  (2) (i) 设过左焦点 \(F\) 的直线 \(l\) 的倾角为 \(\theta\) （斜率大于 0，且 \(Q\) 在第三象限，故 \(0 < \theta < \dfrac{\pi}{2}\)）。\\
  令 \(u = \cos\theta\)，\(v = \sin\theta\)，则有 \(u, v > 0\) 且 \(u^2 + v^2 = 1\)。\\
  以 \(F(-1, 0)\) 为极点，水平向右为极轴建立极坐标系。椭圆极坐标方程的半通径 \(p_0 = \dfrac{b^2}{a} = \dfrac{3}{2}\)，离心率 \(e = \dfrac{1}{2}\)。\\
  由焦点弦极坐标方程可知，从焦点 \(F\) 出发，沿正方向 \((u, v)\) 和反方向 \((-u, -v)\) 到椭圆的极半径距离分别记为：\\
  \(FP = r_1 = \dfrac{p_0}{1 - e\cos\theta} = \dfrac{3}{2 - u}\)；\\
  \(FQ = r_2 = \dfrac{p_0}{1 - e\cos(\theta + \pi)} = \dfrac{3}{2 + u}\)。\\
  因此，\(P = F + r_1(u, v) = (-1 + r_1 u, r_1 v)\)，\(Q = F - r_2(u, v) = (-1 - r_2 u, -r_2 v)\)。\\
  因为点 \(R\) 为 \(PO\) 的延长线与椭圆的交点，由中心对称性得 \(R = -P\)。\\
  计算面积：\\
  \(\triangle PFO\) 以 \(FO = 1\) 为底，点 \(P\) 的纵坐标为高，所以面积为：\\
  \(S_{\triangle PFO} = \dfrac{1}{2} \cdot FO \cdot y_P = \dfrac{1}{2} r_1 v\)。\\
  \(\triangle PQR\) 的面积（以 \(R = -P\) 展开行列式）为：\\
  \(S_{\triangle PQR} = \dfrac{1}{2} |\det(\vec{PQ}, \vec{PR})| = \dfrac{1}{2} |\det(Q - P, -2P)| = |\det(P, Q)| = (r_1 + r_2)v\)。\\
  由题意 \(S_{\triangle PQR} = 3 S_{\triangle PFO}\)，可得：\\
  \((r_1 + r_2)v = \dfrac{3}{2} r_1 v\)。\\
  因为 \(v = \sin\theta > 0\)，消去 \(v\) 得：\\
  \(r_1 + r_2 = \dfrac{3}{2} r_1 \implies r_2 = \dfrac{1}{2} r_1 \implies \dfrac{r_2}{r_1} = \dfrac{1}{2}\)。\\
  代入 \(r_1, r_2\) 的代数式：\\
  \(\dfrac{2 - u}{2 + u} = \dfrac{1}{2} \implies 4 - 2u = 2 + u \implies 3u = 2 \implies u = \dfrac{2}{3}\)。\\
  由此得斜率为：\\
  \(k = \tan\theta = \dfrac{\sqrt{1 - u^2}}{u} = \dfrac{\sqrt{1 - 4/9}}{2/3} = \dfrac{\sqrt{5}}{2}\)。\\
  又直线过 \(F(-1, 0)\)，所以直线 \(l\) 的方程为：\\
  \(y = \dfrac{\sqrt{5}}{2}(x + 1)\)。\\
  (ii) 向量 \(\vec{QP} = P - Q = (r_1 + r_2)(u, v)\)；\\
  向量 \(\vec{QR} = R - Q = -P - Q = \left(2 - (r_1 - r_2)u, -(r_1 - r_2)v\right)\)。\\
  设 \(\angle PQR = \varphi\)，由平面向量的夹角正切公式（外积除以点积）：\\
  \(\tan\varphi = \dfrac{|\det(\vec{QP}, \vec{QR})|}{\vec{QP} \cdot \vec{QR}}\)。\\
  计算分子（行列式绝对值）：\\
  \(|\det(\vec{QP}, \vec{QR})| = (r_1 + r_2) \left| \det\left((u,v), (2,0) - (r_1 - r_2)(u,v)\right) \right| = (r_1 + r_2) \left| \det\left((u,v), (2,0)\right) \right| = 2(r_1 + r_2)v\)。\\
  计算分母（点积）：\\
  \(\vec{QP} \cdot \vec{QR} = (r_1 + r_2) (u, v) \cdot \left[(2, 0) - (r_1 - r_2)(u,v)\right] = (r_1 + r_2) \left[ 2u - (r_1 - r_2)(u^2+v^2) \right] = (r_1 + r_2) [ 2u - (r_1 - r_2) ]\)。\\
  代入正切公式，消去 \((r_1 + r_2)\) 得：\\
  \(\tan\varphi = \dfrac{2v}{2u - (r_1 - r_2)}\)。\\
  由极半径代数式计算差值：\\
  \(r_1 - r_2 = \dfrac{3}{2 - u} - \dfrac{3}{2 + u} = \dfrac{6u}{4 - u^2}\)。\\
  所以分母为：\\
  \(2u - (r_1 - r_2) = 2u - \dfrac{6u}{4 - u^2} = \dfrac{8u - 2u^3 - 6u}{4 - u^2} = \dfrac{2u(1 - u^2)}{4 - u^2}\)。\\
  得：\\
  \(\tan\varphi = \dfrac{2v}{\dfrac{2u(1 - u^2)}{4 - u^2}} = \dfrac{v(4 - u^2)}{u(1 - u^2)} = \dfrac{\sqrt{1 - u^2}(4 - u^2)}{u(1 - u^2)} = \dfrac{4 - u^2}{u\sqrt{1 - u^2}}\)。\\
  令 \(t = u^2 \ (0 < t < 1)\)，则有：\\
  \(\tan^2\varphi = \dfrac{(4 - t)^2}{t(1 - t)}\)。\\
  设函数 \(G(t) = \dfrac{(4 - t)^2}{t(1 - t)} \ (0 < t < 1)\)，对其求导得：\\
  \(G'(t) = \dfrac{2(t-4) \cdot t(1-t) - (4-t)^2(1-2t)}{[t(1-t)]^2} = \dfrac{(t-4)(4 - 7t)}{t^2(1-t)^2}\)。\\
  令 \(G'(t) = 0 \implies t = \dfrac{4}{7}\)。\\
  当 \(t \in \left(0, \dfrac{4}{7}\right)\) 时，\(G'(t) < 0\)，\(G(t)\) 单调递减；\\
  当 \(t \in \left(\dfrac{4}{7}, 1\right)\) 时，\(G'(t) > 0\)，\(G(t)\) 单调递增。\\
  所以当 \(t = \dfrac{4}{7}\) 时，\(G(t)\) 取得最小值，最小值为：\\
  \(G\left(\dfrac{4}{7}\right) = \dfrac{\left(4 - \dfrac{4}{7}\right)^2}{\dfrac{4}{7} \cdot \left(1 - \dfrac{4}{7}\right)} = \dfrac{\left(\dfrac{24}{7}\right)^2}{\dfrac{12}{49}} = \dfrac{576}{12} = 48\)。\\
  所以 \(\tan\angle PQR\) 的最小值为 \(\sqrt{48} = 4\sqrt{3}\)。%
}

\newcommand{\qEllipsePqrTangentDiagnosis}{%
  学生在解决解析几何第二问时，如果直接联立直线与椭圆方程设坐标，硬求三角形面积和夹角，会产生极度庞大的代数计算量，容易因算错或算不下去而放弃。本题利用焦点弦极坐标参数化，配合行列式与向量外积技巧，可将计算量极大缩减。%
}

\newcommand{\qEllipsePqrTangentTeachingNotes}{%
  极坐标法（焦点坐标化）是解决椭圆/双曲线中过焦点弦相关面积及夹角最值问题的顶级神兵。在讲评此题时，应向学生展示这一解析技巧，重点训练如何将倾角 \(\theta\) 表示为 \((u, v)\)，以及利用行列式计算面积和向量外积公式求解夹角正切值。%
}
